\documentclass[11pt]{article}
\usepackage[margin=1in]{geometry}
\usepackage{amsmath,amssymb,amsthm}
\usepackage{hyperref}
\usepackage{xurl}

\newtheorem{theorem}{Theorem}
\newtheorem{lemma}[theorem]{Lemma}
\theoremstyle{definition}
\newtheorem{definition}[theorem]{Definition}
\theoremstyle{remark}
\newtheorem{remark}[theorem]{Remark}

\let\oldus\_
\renewcommand{\_}{\oldus\allowbreak}

\title{A proof of Conjecture 00000000027\\[2pt]
\large $r_3(N) = N\exp(-O(\sqrt{\log N}))$ (Behrend's bound)}
\date{}

\begin{document}
\sloppy
\maketitle

\section{The conjecture and its reading}

Conjecture 00000000027 reads: \emph{$r_3(N)$ (the maximal size of a subset of $[N]$ with no
three-term arithmetic progression) satisfies $r_3(N) = N\cdot\exp(-O(\sqrt{\log N}))$.} The Chinese
text states the same formula.

\paragraph{Reading.} With the usual meaning of the $O$-notation inside a formula, the statement asserts
that there is a function $g$ with $g(N) = O(\sqrt{\log N})$ as $N\to\infty$ and
$r_3(N) = N e^{-g(N)}$ for all large $N$. Because $1 \le r_3(N) \le N$ for $N \ge 1$, this is the same as
the existence of a constant $C$ with $N e^{-C\sqrt{\log N}} \le r_3(N)$ for all large $N$, i.e.\ Behrend's
lower bound. We prove it with the explicit constant $C = 4$, for \emph{every} $N\ge 1$, and with
$0 \le g(N) \le 4\sqrt{\log N}$ (so $g$ is non-negative and the $O$-bound holds with constant $4$).

\paragraph{What is not claimed.} The one-sided notation $-O(\cdot)$ does not assert a matching upper
bound, and we do \emph{not} prove one: the sharpness statement $r_3(N) \le N e^{-c\sqrt{\log N}}$ for some
$c>0$ (a $\Theta$-version) is not claimed here. For context, the Wikipedia article on Salem--Spencer sets
(retrieved; quoted in the attempt file) records upper bounds of the much weaker shape
$\exp(-c(\log N)^{1/12})N$ (Kelley and Meka) with the exponent later improved to $1/9$ (Bloom and Sisask).

\paragraph{Conventions.} $\log$ is the natural logarithm and $\sqrt{\cdot}$ the non-negative square
root. $[N] = \{1,\dots,N\}$. A three-term arithmetic progression is $a, a+d, a+2d$ with common difference
$d \ge 1$ (trivial progressions $a,a,a$ are excluded, as usual). We also show that the value of $r_3(N)$ is
the same for the convention $[N] = \{0,\dots,N-1\}$.

\section{Definitions (as in the Lean file)}

\begin{definition}[\texttt{C27.HasThreeAP}, \texttt{C27.r3}]
A finite set $A\subseteq\mathbb N$ \emph{has a three-term AP} if there are $a, d \in \mathbb N$ with
$d \ge 1$ and $a, a+d, a+2d \in A$. Then
\[ r_3(N) = \max\{\, |A| : A \subseteq \{1,\dots,N\},\ A \text{ has no three-term AP} \,\}, \]
formalized as the supremum of \texttt{card} over the filtered powerset of \texttt{Icc 1 N} (the empty set
qualifies, so this is a genuine maximum).
\end{definition}

Mathlib's \texttt{ThreeAPFree s} means: for all $a,b,c\in s$, $a + c = b + b$ implies $a = b$.
Mathlib's \texttt{addRothNumber s} is the largest size of a \texttt{ThreeAPFree} subset of the finset $s$,
and \texttt{rothNumberNat N = addRothNumber (range N)}, the $\{0,\dots,N-1\}$ version.

\section{Proof}

\begin{lemma}[\texttt{not\_hasThreeAP\_iff}]
For a finset $A\subseteq\mathbb N$: $A$ has no three-term AP iff \texttt{ThreeAPFree A}.
\end{lemma}
\begin{proof}
If $a,a+d,a+2d\in A$ with $d\ge1$, then $a + (a+2d) = (a+d)+(a+d)$ but $a \ne a+d$. Conversely, let
$a,b,c\in A$ with $a+c = 2b$ and $a\ne b$. If $a<b$, put $d = b-a$; then $b = a+d$ and $c = a+2d$. If
$a > b$, then $c < b$; put $d = b - c$, so $b = c+d$ and $a = c + 2d$.
\end{proof}

\begin{lemma}[\texttt{r3\_eq\_addRothNumber}, \texttt{r3\_eq\_rothNumberNat}]
$r_3(N) = \texttt{addRothNumber}(\{1,\dots,N\}) = \texttt{rothNumberNat}(N)$.
\end{lemma}
\begin{proof}
The first equality follows from the previous lemma by comparing the two maxima in both directions. For
the second, $\{1,\dots,N\} = [1, N+1)$ and Mathlib's \texttt{addRothNumber\_Ico} gives
$\texttt{addRothNumber}([a,b)) = \texttt{rothNumberNat}(b-a)$ (translation invariance).
\end{proof}

\begin{lemma}[\texttt{r3\_le}, \texttt{behrend\_r3}, \texttt{r3\_two\_sided}]
For every $N$: $N e^{-4\sqrt{\log N}} \le r_3(N) \le N$.
\end{lemma}
\begin{proof}
The upper bound is Mathlib's \texttt{rothNumberNat\_le}. The lower bound is Mathlib's
\texttt{Behrend.roth\_lower\_bound : (N : R) * exp (-4 * sqrt (log N)) <= rothNumberNat N}, a formalization
(by Y.~Dillies and B.~Mehta, file \texttt{Mathlib/Combinatorics/Additive/AP/Three/Behrend.lean}) of
Behrend's construction [Behrend 1946], transported by the previous lemma.
\end{proof}

\begin{theorem}[\texttt{C27.conjecture\_27}]
Let $g(N) = \log(N / r_3(N))$. Then $g = O(\sqrt{\log N})$ as $N\to\infty$; for every $N \ge 1$,
$0 \le g(N) \le 4\sqrt{\log N}$; and for every $N\ge1$, $r_3(N) = N e^{-g(N)}$.
\end{theorem}
\begin{proof}
For $N \ge 1$ the previous lemma gives $r_3(N) \ge N e^{-4\sqrt{\log N}} > 0$, so $g(N)$ is the logarithm
of a positive number and $N e^{-g(N)} = N \cdot r_3(N)/N = r_3(N)$ (\texttt{r3\_explicit}). From
$r_3(N)\le N$ we get $N/r_3(N) \ge 1$, so $g(N)\ge 0$; from $N \le r_3(N) e^{4\sqrt{\log N}}$ we get
$g(N) \le 4\sqrt{\log N}$. Hence $|g(N)| \le 4\,|\sqrt{\log N}|$ for all $N \ge 1$, which is the
\texttt{IsBigO} statement \texttt{g =O[atTop] fun N => sqrt (log N)}.
\end{proof}

\section{Lean formalization}

File \texttt{lean/Conjecture27/Basic.lean} (137 lines, only \texttt{import Mathlib}), namespace
\texttt{C27}. Main theorem:
\begin{verbatim}
theorem conjecture_27 :
    exists g : Nat -> Real, (g =O[atTop] fun N : Nat => sqrt (log N)) /\
      (forall N, 1 <= N -> 0 <= g N /\ g N <= 4 * sqrt (log N)) /\
      forall N, 1 <= N -> (r3 N : Real) = N * exp (-g N)
\end{verbatim}
(written here in ASCII; the Lean source uses the usual Unicode symbols). Further theorems:
\texttt{not\_hasThreeAP\_iff}, \texttt{r3\_eq\_addRothNumber}, \texttt{r3\_eq\_rothNumberNat},
\texttt{r3\_le}, \texttt{behrend\_r3}, \texttt{r3\_pos}, \texttt{r3\_explicit}, \texttt{r3\_two\_sided}.

\paragraph{Axiom audit.} \texttt{lake build} succeeds, and \texttt{\#print axioms} for
\texttt{C27.conjecture\_27}, \texttt{C27.r3\_two\_sided} and \texttt{C27.r3\_eq\_rothNumberNat} reports
only \texttt{[propext, Classical.choice, Quot.sound]}. There is no \texttt{sorry}, no
\texttt{native\_decide} and no added axiom.

\section*{Attribution and sources}

The result is Behrend's theorem; the contribution here is the faithful statement of the conjecture's
objects and its derivation from Mathlib.
\begin{itemize}
\item F.~A.~Behrend, \emph{On Sets of Integers Which Contain No Three Terms in Arithmetical
Progression}, Proc.\ Natl.\ Acad.\ Sci.\ USA (1946), \url{https://doi.org/10.1073/pnas.32.12.331}
(bibliographic record retrieved from Crossref).
\item Wikipedia, \emph{Salem--Spencer set}, \url{https://en.wikipedia.org/wiki/Salem\%E2\%80\%93Spencer_set}:
``The construction of Salem and Spencer was improved by Felix Behrend in 1946, who found sets of size
$n/e^{O(\sqrt{\log n})}$.''
\item Mathlib, \texttt{Mathlib/Combinatorics/Additive/AP/Three/Behrend.lean} (\texttt{Behrend.roth\_lower\_bound})
and \texttt{Mathlib/Combinatorics/Additive/AP/Three/Defs.lean} (\texttt{rothNumberNat}, \texttt{addRothNumber\_Ico}).
\end{itemize}

\end{document}
